\chapter{Interrogazioni di tipi insiemistico}
\section{Operatori insiemistici}
Permettono di costruire query concatenando due query SQL
\begin{itemize}
    \item UNION
    \item INTERSECT
    \item EXCEPT (Differenza)
\end{itemize}

\subsection{Sintassi}
\begin{lstlisting}[language=sql_generic]
SelctSQL{< UNION / INTERSECT / EXCEPT > [ALL] SelectSQL}
\end{lstlisting}

I duplicati vengono eliminati (a meno che si usi ALL)


\section{UNION}
\subsection{Cosa fa?}
Esegue l'unione dei risultati delle 2 query.

L'operatore UNION impone alcune restrizioni sulle interrogazioni su cui opera:
\begin{itemize}
    \item Le interrogazioni devono restituire lo stesso numero di colonne, e le colonne corrispondenti devono avere lo stesso dominio o domini compatibili
    \item La corrispondenza si basa sulla posizione delle colonne indipendentemente dal loro nome
    \item Se si usa la clausola di ORBER BY queste deve essere usata una sola volta alla fine dell'interrogazione e non alla fine di ogni select
\end{itemize}

\subsection{Sintassi}
\begin{lstlisting}[language=sql_generic]
SelctSQL{ UNION [ALL] SelectSQL}
\end{lstlisting}

\subsection{Esempio}
\begin{lstlisting}[language=sql_generic]
SELECT CodOrd 
FROM Ordine 
WHERE Importo > 500

UNION 

SELECT CodOrd 
FROM Dettaglio 
WHERE Qta > 1000;
\end{lstlisting}

\section{INTERSECT}
\subsection{Cosa fa?}
Esegue l'intersezione dei risultati delle 2 query.

L'operatore INTERSECT impone alcune restrizioni sulle interrogazioni su cui opera:
\begin{itemize}
    \item Le interrogazioni devono restituire lo stesso numero di colonne, e le colonne corrispondenti devono avere lo stesso dominio o domini compatibili
    \item La corrispondenza si basa sulla posizione delle colonne indipendentemente dal loro nome
    \item Se si usa la clausola di ORBER BY queste deve essere usata una sola volta alla fine dell'interrogazione e non alla fine di ogni select
\end{itemize}

\subsection{Sintassi}
\begin{lstlisting}[language=sql_generic]
SelctSQL{ INTERSECT [ALL] SelectSQL}
\end{lstlisting}

\subsection{Esempio}
\begin{lstlisting}[language=sql_generic]
SELECT CodOrd 
FROM Ordine 
WHERE Importo > 500

INTERSECT

SELECT CodOrd 
FROM Dettaglio 
WHERE Qta > 1000;
\end{lstlisting}

\section{EXCEPT}
\subsection{Cosa fa?}
Esegue la differenza dei risultati delle 2 query.

L'operatore EXCPET impone alcune restrizioni sulle interrogazioni su cui opera:
\begin{itemize}
    \item Le interrogazioni devono restituire lo stesso numero di colonne, e le colonne corrispondenti devono avere lo stesso dominio o domini compatibili
    \item La corrispondenza si basa sulla posizione delle colonne indipendentemente dal loro nome
    \item Se si usa la clausola di ORBER BY queste deve essere usata una sola volta alla fine dell'interrogazione e non alla fine di ogni select
\end{itemize}

\subsection{Sintassi}
\begin{lstlisting}[language=sql_generic]
SelectSQL { EXCEPT [ ALL ] SelectSQL }
\end{lstlisting}

\subsection{Esempio}
\begin{lstlisting}[language=sql_generic]
SELECT Nome 
FROM Impiegato

EXCEPT 

SELECT Cognome AS Nome 
FROM Impiegato;
\end{lstlisting}